\documentclass[10pt]{amsart}
\usepackage{calc,amssymb,amsmath,ulem,stmaryrd,fullpage}
\usepackage{enumerate}
\usepackage{alltt}
\usepackage{multicol}
\RequirePackage[dvipsnames,usenames]{color}
\let\mffont=\sf
\normalem
%\input{kmacros3.sty}
\input{xy}
\xyoption{all}
\usepackage{hyperref}
\usepackage{graphicx}
\usepackage[all,cmtip]{xy}
\usepackage{verbatim}
\usepackage[left=0.95in,top=0.95in,right=0.95in,bottom=0.95in]{geometry}
\newcommand{\tauCohomology}{T}
\newcommand{\FCohomology}{S}


\theoremstyle{theorem}
%\newtheorem*{mainthm*}{Main Theorem}

\newcommand{\note}[1]{\marginpar{\sffamily\tiny #1}}

\def\todo#1{\textcolor{Mahogany}%
{\sffamily\footnotesize\newline{\color{Mahogany}\fbox{\parbox{\textwidth-15pt}{#1}}}\newline}}

\renewcommand{\O}{\mathcal O}
\newcommand{\ie}{{\itshape i.e.} }
\newcommand{\roundup}[1]{\lceil #1 \rceil}
\newcommand{\rounddown}[1]{\lfloor #1 \rfloor}
\renewcommand{\to}[1][]{\xrightarrow{\ #1\ }}
\newcommand{\onto}[1][]{\protect{\xrightarrow{\ #1\ }\hspace{-0.8em}\rightarrow}}
\newcommand{\into}[1][]{\lhook \joinrel \xrightarrow{\ #1\ }}
%\newcommand{DBDual}[1]{{\underline{\omega}_{#1}}}
\newcommand{\DBDual}[1]{{{\uuline \omega}{}^{\mydot}_{#1}}}
\newcommand{\Tt}{{\mathfrak{T}}}
%\newcommand{\bn}{{\mathfrak{n}} }
\newcommand{\sol}[1]{ \vskip 6pt{\color{blue} {\bf Solution:} #1} \vskip 6pt}

\newcommand{\bR}{{\mathbb{R}}}
\newcommand{\va}{{\vec{a}}}
\newcommand{\vb}{{\vec{b}}}
\newcommand{\vzero}{{\vec{0}}}
\newcommand{\vi}{{\vec{i}}}
\newcommand{\vj}{{\vec{j}}}
\newcommand{\vk}{{\vec{k}}}
\newcommand{\vh}{{\vec{h}}}
\newcommand{\vr}{{\vec{r}}}
\newcommand{\vu}{{\vec{u}}}
\newcommand{\vv}{{\vec{v}}}
\newcommand{\vw}{{\vec{w}}}
\newcommand{\vx}{{\vec{x}}}
\newcommand{\vy}{{\vec{y}}}
\newcommand{\vz}{{\vec{z}}}
\newcommand{\vT}{{\vec{T}}}
\newcommand{\vN}{{\vec{N}}}
\newcommand{\vF}{{\vec{F}}}
\newcommand{\vG}{{\vec{G}}}
\newcommand{\bF}{\mathbb{F}}
\newcommand{\bQ}{\mathbb{Q}}
\newcommand{\bZ}{\mathbb{Z}}
\newcommand{\bC}{\mathbb{C}}


\begin{document}
\title{Homework \#2 -- Math 5320\\Spring 2020}
\author{Due Friday, February 14th, in Gradescope}

\maketitle


\noindent
{\bf 1.}  Suppose that $k$ is any field.  For any polynomial $f \in k[x]$, $f = \sum_{k =0}^n a_k x^k$ we can define the \emph{formal derivative} $f' = \sum_{k = 1}^n k a_k x^{k-1}$, which is the usual derivative formula when $k = \bR$.  
\vskip 3pt
{\bf (a)}  Prove the product rule $(fg)' = f'g + fg'$, and 
\vskip 9pt
{\bf (b)}  prove the chain rule $(f(g(x)))' = f'(g(x)) \cdot g'(x)$.
\newpage
\noindent
{\bf 2.}  Suppose $k$ is a field, $\alpha \in k$, $f \in k[x]$ and $f(\alpha) = 0$.  Prove that $\alpha$ is a multiple root of $f$ if and only if $f'(\alpha) = 0$.  
\newpage
\noindent
{\bf 3.}  Classify all rings with 4 elements up to isomorphism.
\newpage
\noindent
{\bf 4.}  Let $S = \bC[x,y]$ and suppose $I = (y^2 + x^3 - 17)$.  Which of the following are maximal ideals in $S/I$.
\vskip 3pt
{\bf (a)}  $(x-1,y-4)$
\vskip 3pt
{\bf (b)}  $(x+1, y+4)$
\vskip 3pt
{\bf (c)}  $(x^3 - 17, y^2)$
\newpage
\noindent
{\bf 5.}  Fix $a,b, c \in \bC$ (with either $a$ or $b$ nonzero) and fix $0 \neq f(x,y) \in \bC[x,y]$.  Let $L$ be the complex ``line'' $\{(x,y) \in \bC^2 \;|\; ax + by + c = 0\}$ and let $C$ be the complex ``curve'' $\{(x,y) \in \bC^2 \;|\; f(x,y) = 0\}$.  Prove that either $L$ and $C$ intersect at $\leq \deg f$ points or at infinitely points (in which case $L \subseteq C$).
\newpage
\noindent
{\bf 6.}  Which maximal ideals in the polynomial ring $\bC[x,y]$ contain both $x^2 + y^2 -5$ and $xy - 2$?
\newpage
\noindent
{\bf 7.}  Euclid proved that there are infinitely many prime integers $n$ as follows.  If $p_1, \dots, p_n$ are all the prime integers, then any prime factor of $p_1 \cdot p_2 \cdots p_n + 1$ is different from all the $p$.  Use this strategy to prove that if $k$ is a field, then $k[x]$ has infinitely many monic irreducible polynomials.
\newpage
\noindent
{\bf 8.}  Prove that $\bZ[\sqrt{-2}]$ is a Euclidean domain. 


\end{document} 